\documentclass[tikz,border=5mm]{standalone}
\usetikzlibrary{angles,quotes}

\begin{document}
\begin{tikzpicture}[scale=1.5, >=stealth]

  % 1. Definition der Eckpunkte (Allgemeines Dreieck)
  \coordinate[label=below left:$A$] (A) at (0,0);
  \coordinate[label=below right:$B$] (B) at (5,0);
  \coordinate[label=above:$C$] (C) at (1.5,3.5);

  % 2. Zeichnen der Dreiecksseiten und deren Beschriftung
  \draw[thick] (A) -- (B) node[midway, below=2pt] {$c$}
                    -- (C) node[midway, above right=1pt] {$a$}
                    -- cycle node[midway, above left=1pt] {$b$};

  % 3. Zeichnen und Beschriften der Innenwinkel
  % Die 'angles'-Bibliothek zeichnet automatisch Bögen gegen den Uhrzeigersinn.
  % 'angle eccentricity' bestimmt den Abstand der Beschriftung vom Scheitelpunkt.
  \pic[draw, "$\alpha$", angle radius=0.7cm, angle eccentricity=1.4] {angle = B--A--C};
  \pic[draw, "$\beta$", angle radius=0.7cm, angle eccentricity=1.4] {angle = C--B--A};
  \pic[draw, "$\gamma$", angle radius=0.6cm, angle eccentricity=1.4] {angle = A--C--B};

\end{tikzpicture}
\end{document}